\documentclass[11pt]{article}
\usepackage[margin=1in]{geometry}
\usepackage{amsmath,amssymb,amsthm,booktabs}
\newtheorem{lemma}{Lemma}
\newtheorem{proposition}{Proposition}
\newtheorem{theorem}{Theorem}
\newtheorem{conjecture}{Conjecture}
\title{Access Schedules and Stash Growth in the Current Lane ORAM}
\author{}
\date{July 2026}

\begin{document}
\maketitle

\paragraph{Model.}
There are $N$ initialized logical keys.  On every access, the requested key is
read on its old leaf path and assigned a fresh independent uniform leaf label.
The logical schedule is fixed independently of these random labels.  This is
the threat model used for the empirical experiment; an adaptive schedule that
can observe secret labels is outside the ORAM model.

\begin{lemma}[Uniform physical paths]
Let $(a_t)_{t\geq 1}$ be a fixed logical schedule independent of the ORAM
coins.  For each key $a$, let
$U_{a,0},U_{a,1},\ldots$ be independent uniform elements of $[N]$, where
$U_{a,0}$ is its initial label and $U_{a,j}$ is the label sampled by its $j$th
access.  If
\[
  r_t=\left|\{s<t:a_s=a_t\}\right|,
\]
then the path read at time $t$ is $X_t=U_{a_t,r_t}$ and its fresh label is
$Y_t=U_{a_t,r_t+1}$.  Consequently, $(X_t)_{t\geq1}$ is i.i.d. uniform and
$X_t$ is independent of $Y_t$.
\end{lemma}

\begin{proof}
The pairs $(a_t,r_t)$ are distinct: successive occurrences of one key have
successive occurrence numbers, while different keys have different first
coordinates.  Thus the variables $X_t$ select distinct members of the
independent family $(U_{a,j})$.  The two variables defining $X_t$ and $Y_t$
are also distinct.
\end{proof}

This lemma does \emph{not} make the entire Lane ORAM transition sequence
schedule-independent.  If $t'$ is the next access to $a_t$, then
$X_{t'}=Y_t$.  The reuse distance $t'-t$ therefore controls how long a block
with label $Y_t$ is exposed to intervening lane-local evictions.

\begin{lemma}[Repeated permutations are equivalent]
Under exchangeable initialization and independent uniform labels, any two
logical schedules obtained by endlessly repeating a permutation of all $N$
keys have the same stash distribution.
\end{lemma}

\begin{proof}
Rename each key according to the bijection taking the first permutation to the
second.  The initial labels and all per-key label streams have a distribution
invariant under this renaming.  The current Lane ORAM transition examines a
key only for equality and otherwise uses its label, so the renamed execution
has exactly the second schedule and the same stash occupancies.
\end{proof}

\begin{proposition}[Maximally nonrepeating schedules]
Suppose an infinite schedule over $N$ keys has no repeated key in any window of
$N$ consecutive accesses.  Then it is a repeated permutation: every length
$N$ window contains all keys and $a_{t+N}=a_t$ for every $t$.
\end{proposition}

\begin{proof}
A length-$N$ window contains $N$ distinct elements of an alphabet of size
$N$, hence contains every key.  The windows beginning at $t$ and $t+1$ differ
only by removing $a_t$ and adding $a_{t+N}$; equality of their underlying sets
forces $a_{t+N}=a_t$.
\end{proof}

Thus round-robin has the largest possible minimum reuse distance.  It also
represents every repeated permutation schedule.  The experiments confirm that
it is strongly adverse: at $N=4096$, $B=2$, $Z=3$, mean stash occupancy was
about $1885$ for four different repeated permutations, $1660$ for i.i.d.
uniform logical keys, $0.30$ for a uniformly accessed hot set of 16 keys, and
$0$ for one repeated key.

\paragraph{What is not proved.}
These facts do not prove that round-robin is the global worst-case schedule for
Lane ORAM.  The duplicate-deletion argument cited for Path and Circuit ORAM
uses an infinite-bucket representation followed by order-independent global
greedy postprocessing.  In the current Lane ORAM, a block keeps a persistent
lane in the tree, root admission takes the first stash block, and holes move
only within one lane.  Deleting a duplicate access can therefore change both
future lane assignments and hole chains.  No monotone coupling between reuse
distance and Lane stash occupancy is presently known.  The defensible claim is
that round-robin is \emph{empirically adversarial} and maximally nonrepeating,
not that it is a proved global maximizer.

\paragraph{A rigorous initialization lower bound.}
The current API loads a previously absent key by calling
\texttt{update(0,new\_pos,key)}.  Loading distinct keys therefore reads path
zero every time until the workload begins.

\begin{theorem}[Sequential-loading failure]
Load $N$ distinct absent keys, never revisiting a key during loading, and give
each one a fresh independent uniform label.  Let $A$ be the number of blocks
admitted from the stash into the tree and let $Q=N-A$ be the post-loading stash
occupancy.  Put $p=(B-1)/B$.  Then
\[
  \mathbb{E}[A]\leq\frac{Z}{p}=\frac{ZB}{B-1},
  \qquad
  \mathbb{E}[Q]\geq N-\frac{ZB}{B-1},
\]
and, for every integer $m\geq0$,
\[
  \Pr[Q<N-m]=\Pr[A>m]
  \leq \Pr[\operatorname{Bin}(m,p)<Z].
\]
Consequently, the probability that a stash of capacity $S$ survives loading is
at most
\[
  \Pr[\operatorname{Bin}(N-S-1,p)<Z].
\]
\end{theorem}

\begin{proof}
Consider a block admitted into a free root lane.  With probability $p$, the
first base-$B$ digit of its fresh label is nonzero.  Such a block cannot move
below the root on path zero: its common prefix with zero has length less than
$\log_2B$.  It is not requested again during loading, so that root lane is
permanently sealed.  The labels of successively admitted blocks are independent.
After $Z$ sealing labels, all root lanes are sealed and no further block can be
admitted.  Any blockage deeper in a lane can only stop admissions earlier.
Thus $A$ is stochastically dominated by the number of Bernoulli-$p$ trials
needed to obtain $Z$ successes.  Its expectation is at most $Z/p$, and
$A>m$ implies fewer than $Z$ successes among the first $m$ trials, proving the
tail bound.  During loading no tree block is removed, so stash occupancy is
nondecreasing and a capacity-$S$ execution survives only if $A\geq N-S$.
\end{proof}

For $N=4096$ and $S=20$, the theorem bounds survival by $2^{-4052}$ for
$(B,Z)=(2,3)$; by $2^{-8100}$ and $2^{-8089}$ for $(4,5)$ and $(4,6)$; and by
$2^{-12122}$ and $2^{-12110}$ for $(8,9)$ and $(8,10)$.  These are upper bounds
on survival, not estimates.  The effectively unbounded-stash stationary
experiment below is therefore hypothetical: the real $S=20$ instance almost
certainly fails before reaching its warm-up phase.

This lower bound persists for an initial portion of a later cyclic workload.
One update can reduce post-writeback stash occupancy by at most $Z$, because
removal and reinsertion do not reduce the stash before root admission and at
most $Z$ root slots can drain it.  Therefore, deterministically,
$Q_w\geq Q_N-Zw$ after $w$ further updates.  In particular, the stash remains
linear for the first $\Theta(N)$ operations.

\paragraph{Structural security evidence.}
After initialization, block conservation gives the exact identity
\[
  S_t=N-\sum_{\ell=1}^{Z}T_{t,\ell},
\]
where $S_t$ is post-eviction stash occupancy and $T_{t,\ell}$ is the number of
blocks resident in tree lane $\ell$.  The five trees have between $6N$ and
$11.43N$ physical block slots, so the observed large stash is not ordinary
capacity overload.  It is a failure of the fixed-lane admission and eviction
dynamics to exploit those slots.

The size-scaling experiments give the following mean stash fractions:
\begin{center}
\begin{tabular}{lcc}
\toprule
Configuration & largest exact $N$ tested & $\mathbb{E}[S]/N$ \\
\midrule
$B=2,Z=3$   & $2^{14}$ & $0.501$ \\
$B=4,Z=5$   & $2^{14}$ & $0.740$ \\
$B=4,Z=6$   & $2^{14}$ & $0.689$ \\
$B=8,Z=9$   & $2^{15}$ & $0.833$ \\
$B=8,Z=10$  & $2^{15}$ & $0.813$ \\
\bottomrule
\end{tabular}
\end{center}
This directly contradicts the $N$-independent stash behavior sought from the
Circuit ORAM methodology.

\begin{conjecture}[Empirical fixed-lane scaling]
Let $k=\log_2B$ and $L=\log_BN$.  Under repeated-permutation access, the mean
tree occupancy of one lane satisfies
\[
  \frac{\mathbb{E}[T_{t,\ell}]}{N}
  =\Theta\!\left(L^{-k/2}\right).
\]
Consequently, for fixed $B$ and $Z$,
\[
  \frac{\mathbb{E}[S_t]}{N}
  =1-\Theta\!\left(\frac{Z}{L^{k/2}}\right)\longrightarrow1.
\]
\end{conjecture}
The measured constants for $B=2,4,8$ are approximately $0.62$, $0.36$, and
$0.21$, respectively.  This law fits all tested sizes and independent one-lane
diagnostics, but it is an empirical conjecture rather than a security theorem.
A proof would require a stationary analysis of the lane-local hole process and
its root-admission queue.

\paragraph{Failure probability.}
The implementation fails at the transient insertion point, before root
draining.  Therefore a stash of capacity $R$ fails on an operation exactly
when the insertion demand exceeds $R$; post-writeback occupancy can understate
this event.  A stationary snapshot estimate
$\delta_R=\Pr[\text{demand}>R]$ also is not an ever-failure probability.  For
$T$ operations, only the general union bound
\[
  \Pr[\text{some failure in }T]\leq T\delta_R
\]
follows without a dependence analysis.  Conversely, the experiments observed
demand above $20$ on every one of $2^{27}$ measured round-robin operations for
every requested configuration at $N=4096$.  Hence the current default
$S=20$ is empirically inadequate by orders of magnitude.

\end{document}
